\section{Method}\label{sec:method}

The system has three layers: a point-forecaster pool (held fixed; not the
contribution), an online regime estimator, and the calibration layer that
is the subject of the paper. Figure~\ref{fig:coverage} previews the
outcome; this section specifies the machinery that produces it.

\subsection{Forecasting layer (fixed)}

The target is next-day log realized variance $y_{i,t+1} = \log
\mathrm{RV}_{i,t+1}$ per stock $i$. The pool combines a per-stock HAR
regression \citep{corsi2009simple} and a pooled gradient-boosted model
(LightGBM) by exponentially weighted aggregation on QLIKE loss
\citep{patton2011volatility}, refit on an expanding window with a hard
train/test date assertion. The pool ties its components (QLIKE 0.191 vs
0.194/0.188; Table~\ref{tab:e0}); everything downstream applies unchanged
to any member (Section~\ref{sec:ablations} conformalizes HAR and LightGBM
alone), to a zero-shot foundation model (Section~\ref{sec:tsfm}), and to
the return stream directly for VaR (Section~\ref{sec:var}).

\subsection{Regime layer}

At the close of day $t$ the market state comprises the trailing-750-day
VIX percentile, the market realized-volatility percentile, and
cross-sectional dispersion, all strictly backward-looking. The canonical
estimator is transparent: $K = 4$ regimes as VIX-percentile bins with
edges $(0.5, 0.8, 0.95)$, aligned with the evaluation slices. A soft
alternative (an online Gaussian HMM refit on an expanding window every
quarter, issuing \emph{filtered} probabilities $\pi_t(k)$) yields the
same conclusions (Section~\ref{sec:oracle}), so transparency costs
nothing. An optional forward-looking \emph{acute} group, active when the
9-day implied volatility exceeds 30-day (VIX9D/VIX $> 1$), sharpens the
stress upper tail (Section~\ref{sec:onset}). The calibrator consumes any
membership vector $\pi_t \in \Delta^{K-1}$; nothing assumes the regimes
are correct, and Section~\ref{sec:theory} states guarantees relative to
the estimated memberships.

\subsection{Calibration layer}

\paragraph{Scores.} Each stock-day produces a standardized nonconformity
score $s_{i,t} = (y_{i,t} - m_{i,t}) / \hat\sigma_{i,t}$, where $m_{i,t}$
is the pool forecast and $\hat\sigma_{i,t}$ a trailing MAD scale of that
stock's own recent scores. Standardization puts all stocks on a comparable
scale, which is what makes cross-sectional pooling meaningful.

\paragraph{Pooled per-regime tracking.} The layer maintains per-regime
upper and lower thresholds $q_t^{\mathrm{hi}}(k), q_t^{\mathrm{lo}}(k)$
and per-stock offsets $\delta_i$, issues stock $i$ the interval
\[
  m_{i,t} \pm \bigl( \textstyle\sum_k \pi_t(k)\, q_t(k) + \delta_i \bigr)
  \hat\sigma_{i,t},
\]
and after observing outcomes takes one gradient step per \emph{stock-day}
(Section~\ref{sec:theory}, eq.~\eqref{eq:update}): every regime's
threshold moves proportionally to its membership and to the day's summed
coverage error across the panel. This per-observation stepping is the
pooling mechanism: on a 100-name panel, the stress regime receives
roughly $15{,}000$ updates over the sample where a single stock would
provide $\sim$150. Offsets absorb residual per-stock idiosyncrasy and are
L2-shrunk toward zero (empirically negligible;
Table~\ref{tab:ablations}).

\paragraph{Self-tuning rates.} Rare regimes need fast learning rates and
extreme quantiles need large steps; hand-tuning those is the kind of
researcher degree of freedom we want to eliminate. Each regime therefore
keeps a bank of expert thresholds, one per rate in a fixed dyadic grid,
combined by exponential weights on pinball loss. This is DtACI
\citep{gibbs2022conformal} run \emph{per regime}, with all expert updates,
weight updates, and uniform mixing gated by $\pi_t(k)$ so that a regime
learns only in its own (membership-weighted) time. Because a
weight-averaged threshold can be dragged by slow experts during shifts, a
per-regime additive corrector updates on the \emph{issued} interval's
misses at one universal rate, the quantile-tracking-plus-integrator
structure of conformal PID control \citep{angelopoulos2023conformal}. The
corrector, not the bank, carries the coverage guarantee
(Remark~\ref{rem:adaptive}); the bank carries adaptivity. The full
configuration has \emph{no tuned parameters}, and its learned effective
rates reproduce what hand-tuning found: stress-regime rates roughly twice
calm rates, and 1\%-tail rates roughly twice 5\%-tail rates
(Section~\ref{sec:results}).

\paragraph{One-sided head and VaR.} The risk application tracks a single
upper threshold at level $\alpha \in \{0.05, 0.01\}$ on standardized
return losses, mapping through the location--scale construction to
next-day VaR. All machinery is otherwise identical
(Proposition~\ref{prop:onesided}).

\paragraph{Reductions.} With $K = 1$ the layer is pooled marginal
tracking; with hard memberships, a single-rate grid, and the corrector
off, it reduces \emph{exactly} to the fixed-rate path (regression-tested
to $10^{-12}$); with pooling off it collapses to per-stock Mondrian ACI.
These reductions are the ablation grid of
Section~\ref{sec:ablations}.
